\section{Bounded excess for saturated complete intersections}
\label{sci:sec}

There is also a tractable higher-codimension saturation problem when the
input ideal contains a regular sequence with bounded excess over the
anticanonical degree. A bounded degree gap gives a deterministic
coefficient algorithm. A finite projection and a deterministic interpolation
grid remove the degree bounds on all other inputs. A search over degree
profiles removes the need to supply the regular sequence; only the
selection of a good projection centre remains randomized in the general
bounded-excess algorithm. The input need not generate the saturated
ideal. A small graded annihilator replaces the unknown regular generators.

Write \(S=\Fp[x_0,\ldots,x_N]\) and \(\mathfrak m=(x_0,\ldots,x_N)\).
Let \(S_d\) be its homogeneous degree-\(d\) part for \(d\geq0\), and
put \(S_d=0\) for \(d<0\). All displayed inputs are explicit homogeneous
sparse polynomials, with exponents and degrees encoded in binary.
On the geometric promises below, every ambient coordinate occurs in
the input list. If one did not, its coordinate vertex would be a cone
vertex of \(X\). At that vertex the nonlinear minimal equations have
zero gradient, so smoothness would force all \(c\) minimal degrees to
be one, contrary to their anticanonical sum \(N+1>c\). Thus \(N+1\)
is bounded by the written sparse length even when coordinate indices
are encoded in binary.

\begin{theorem}[Recovery from a regular sublist with a bounded degree gap]
\label{sci:thm:gap}
Fix a prime \(p\), a codimension \(c\geq2\), and integers \(e,B\geq0\).
Let \(K=(u_1,\ldots,u_m)\), where the \(u_i\) are nonzero homogeneous
sparse forms. Promise that \(K^{\mathrm{sat}}=I_X\), where
\(X\subset\PP^N_{\Fp}\) is a smooth geometrically integral global
complete intersection of codimension \(c\), dimension \(n=N-c\geq1\),
and anticanonical multidegree. Give indices of a sublist
\(f_1,\ldots,f_c\) which is a regular sequence, and suppose
\[
 D:=\sum_{j=1}^c\deg f_j=N+1+e,\qquad
 \deg u_i\leq \min_j\deg f_j+B\quad(1\leq i\leq m).
\]
Then the scalar of Frobenius on \(H^n(X,\OO_X)\) is computable in
deterministic polynomial sparse-bit time. The algorithm also constructs
a homogeneous \(q\in S_e\) with
\[
 q\notin J,\qquad J:q=I_X,\qquad J=(f_1,\ldots,f_c).
\]
It does not recover the minimal generators of \(I_X\).
Regularity of the indexed sublist is a separate promise, not a
consequence of the saturation promise.
\end{theorem}

We first give the algebra and the Frobenius calculation used in the
algorithm. The intermediate complete intersection \(Y=\Proj(S/J)\)
may be singular, reducible, or nonreduced.

\begin{lemma}[The degree-\(e\) annihilator line]
\label{sci:lem:annihilator}
Under the geometric and regularity hypotheses above,
\[
 J:K=J:I_X,\qquad
 \dim_{\Fp}\bigl((J:K)_e/J_e\bigr)=1.
\]
Every \(q\in(J:K)_e\setminus J_e\) satisfies \(J:q=I_X\).
\end{lemma}
\begin{proof}
The quotient \(S/J\) is Cohen--Macaulay of dimension \(n+1\geq2\), so
\(J\) is saturated. If \(aK\subseteq J\) and \(v\in K^{\mathrm{sat}}\),
then \(\mathfrak m^r v\subseteq K\) for some \(r\), whence
\(\mathfrak m^r av\subseteq J\). Saturation gives \(av\in J\).
This proves \(J:K=J:I_X\).

Choose homogeneous regular generators \(g_1,\ldots,g_c\) of \(I_X\),
of degrees \(d_j\) with \(\sum_jd_j=N+1\), only for the proof.
Write \(f=Hg\) with homogeneous matrix entries of degrees
\(\deg f_i-d_j\), using zero for a negative degree, and put
\(\Delta=\det H\). Every nonzero determinant term has degree \(e\).
Wiebe's annihilator identities give
\[
 J:I_X=J+(\Delta),\qquad J:\Delta=I_X.
\]
Here one may use \citet[Theorem~11 and Corollary~12, \S2--3]{sci-wiebe};
their statement applies in an arbitrary ring when \(f\) is regular,
without a locality assumption. In particular \(\Delta\notin J\),
since \(I_X\ne S\). Taking degree \(e\) gives
\((J:I_X)_e/J_e=\Fp\overline{\Delta}\).
Any nonzero class is a ground-field unit times \(\overline{\Delta}\),
and has the same annihilator.
\end{proof}

Write \(\mathbf1=(1,\ldots,1)\in\mathbb Z^{N+1}\).
Define the \(\Fp\)-linear monomial operator
\[
 \mathcal T(x^a)=
 \begin{cases}
 x^{(a-(p-1)\mathbf1)/p},
       &a_i\equiv p-1\pmod p\text{ for every }i,\\
 0,&\text{otherwise}.
 \end{cases}
\]
In the first case each exponent is nonnegative. Thus
\(\mathcal T(h^p v)=h\mathcal T(v)\).

\begin{lemma}[The Frobenius adjoint on an arbitrary complete intersection]
\label{sci:lem:trace}
Let \(f_1,\ldots,f_c\) be a homogeneous regular sequence with
\(d_j=\deg f_j\), let \(\sum_jd_j=N+1+e\) with \(e\geq0\), and
suppose \(N-c\geq1\).
The dual of \(H^{N-c}(Y,\OO_Y)\) is naturally \(S_e/J_e\). Under this
identification the adjoint of Frobenius is
\[
 \mathcal C_Y(\bar h)
 =\overline{\mathcal T\!\left(h(f_1\cdots f_c)^{p-1}\right)}.
 \label{sci:eq:trace}
\]
If \(J\subseteq I_X\) as above, the line
\((J:I_X)_e/J_e\) is invariant under this operator, and its eigenvalue
is the Frobenius scalar \(\alpha\) of \(X\). In particular,
\[
 \mathcal T\!\left(q(f_1\cdots f_c)^{p-1}\right)
       \equiv\alpha q\pmod{J_e}
 \quad\text{for }q\in(J:I_X)_e\setminus J_e.
\]
This is the trace on a Gorenstein dualizing sheaf; smoothness of \(Y\)
and a description by regular differential forms are not required.
\end{lemma}
\begin{proof}
Put \(d_j=\deg f_j\) and \(D=\sum_jd_j\).
The Koszul resolution and the cohomology of line bundles on
\(\PP^N\) identify the top coherent group with
\[
 \ker\!\left(
 H^N(\PP^N,\OO(-D))
 \longrightarrow
 \bigoplus_jH^N(\PP^N,\OO(-D+d_j))
 \right).
 \tag{\(\ast\)}
\]
Indeed the negative line bundles have cohomology only in row \(N\);
the term \(\OO\) has its separate row-zero class. Since \(N-c\geq1\),
no higher differential connects that row-zero term to the displayed
leftmost kernel. The perfect Laurent coefficient pairing identifies
\(H^N(\OO(-D))^\vee\) with \(S_e\). The dual of the displayed map
is \((a_j)\mapsto\sum_j f_ja_j\), with
\(a_j\in S_{e-d_j}\), so the dual of \((\ast)\) is exactly \(S_e/J_e\).

Apply \(p\)th power to the Koszul resolution and compare the resolution
of \((f_1^p,\ldots,f_c^p)\) with that of \(f\).
On its last term the comparison multiplies by
\((f_1\cdots f_c)^{p-1}\). Thus Frobenius on the ambient top
cohomology is
\[
 v\longmapsto (f_1\cdots f_c)^{p-1}v^p,
\]
and restricts to \((\ast)\). Pair a Laurent basis element
\(x^{-b-\mathbf1}\), \(|b|=e\), with a polynomial \(h\in S_e\).
The adjoint coefficient at \(x^b\) is the coefficient at
\(x^{pb+(p-1)\mathbf1}\) in \(h(f_1\cdots f_c)^{p-1}\).
This is precisely \(\mathcal T\). The resulting polynomial has degree
\[
 \frac{e+(p-1)(N+1+e)-(p-1)(N+1)}p=e.
 \]
The formula also descends directly modulo \(J_e\): replacing \(h\)
by \(h+f_j a\) changes the input by
\(f_j^p a\prod_{k\ne j}f_k^{p-1}\), whose image under \(\mathcal T\)
is divisible by \(f_j\). This derivation is the Koszul coefficient
calculation of Proposition~\ref{sp:prop:ci-coefficient}, now retaining
the possible nonzero groups from proper Koszul terms.

For the last assertion, choose \(g\) and \(H\) as in
Lemma~\ref{sci:lem:annihilator}. The map of Koszul resolutions
inducing \(\OO_Y\to\OO_X\) has top term multiplication by
\(\Delta=\det H\). On top ambient cohomology it sends
\(v\) to \(\Delta v\). Consequently the dual map from
\(H^n(X,\OO_X)^\vee=\Fp\) to \(S_e/J_e\) sends \(1\) to
\(\bar\Delta\). This is nonzero by Wiebe's identity.
Restriction \(\OO_Y\to\OO_X\) commutes with Frobenius.
Its dual therefore identifies \(\Fp\bar\Delta\) as an invariant line
with eigenvalue \(\alpha\). The ground field is \(\Fp\), so adjunction
and rescaling the nonzero generator do not change that scalar.
\end{proof}

\begin{proof}[Algorithm and proof of Theorem~\ref{sci:thm:gap}]
Introduce the coefficients of \(q\in S_e\) and, for every \(i,j\),
the coefficients of
\[
 a_{ij}\in S_{e+\deg u_i-\deg f_j}.
 \]
Solve the homogeneous linear coefficient system
\[
 q u_i=\sum_{j=1}^c a_{ij}f_j\qquad(1\leq i\leq m).
 \tag{\(\dagger\)}
\]
Projection of its solution space onto the \(q\)-coordinates is exactly
\((J:K)_e\): homogeneous ideal membership at the displayed degree
has precisely such coefficient forms. Every auxiliary degree is at
most \(e+B\). Its full monomial space therefore has polynomial size
for fixed \(e,B\). Products with sparse input forms generate the
coefficient rows explicitly, with no dense high-degree expansion.

Compute \(J_e=\sum_j f_jS_{e-\deg f_j}\) in the same way.
Lemma~\ref{sci:lem:annihilator} guarantees that the projected solution
space modulo \(J_e\) has dimension exactly one. Choose any projected
solution outside \(J_e\), and call it \(q\).
Next compute
\[
 r=\mathcal T\!\left(q(f_1\cdots f_c)^{p-1}\right)
 \]
by sparse term enumeration, and reduce \(q,r\) modulo the computed
linear subspace \(J_e\). The nonzero vector \(\bar q\) and the vector
\(\bar r\) satisfy \(\bar r=\alpha\bar q\).
One nonzero coordinate of \(\bar q\) gives \(\alpha\), and the other
coordinates verify the same relation.

For an explicit size bound, put
\[
 A=\binom{N+e+B}{e+B},\qquad
 A_0=\binom{N+e}{e}.
 \]
The system has at most \(A_0+mcA\) unknown coefficients and at most
\(A_0T+mA\sum_jt_j\leq(m+1)TA\) coefficient rows before discarding repetitions,
where \(T\) is the total input term count. Dense elimination is
polynomial in these quantities and in the sparse input bit length.
The quotient computation has at most \(A_0\) coordinates. If \(t_j\)
is the term count of \(f_j\), the Frobenius step uses at most
\[
 A_0\prod_{j=1}^c t_j^{p-1}\leq A_0T^{c(p-1)}
 \]
ordered term choices. Exponent additions, congruence tests and
divisions by \(p\) have polynomial bit cost. Thus the full algorithm
is deterministic polynomial time for fixed \(p,c,e,B\).
No verification of the geometric or regularity promises is asserted.
\end{proof}

\begin{corollary}[An explicit colon certificate without a degree gap]
\label{sci:cor:certificate}
Fix \(p,c,e\). In the same geometric setting, give a homogeneous sparse
regular sequence \(f_1,\ldots,f_c\) with degree sum \(N+1+e\), and an
explicit homogeneous \(q\in S_e\) promised to satisfy
\[
 (f_1,\ldots,f_c):q=I_X.
 \]
Then the Frobenius scalar of \(X\) is computable in deterministic
polynomial sparse-bit time. Any additional equations whose saturation
is \(I_X\) may have arbitrary binary-encoded degrees; they are not used
in the scalar calculation. The condition on \(q\) is an annihilator
condition modulo \(J\), not the condition \(q|_X\ne0\).
\end{corollary}
\begin{proof}
The colon condition implies \(J\subseteq I_X\) and
\(q\in(J:I_X)_e\setminus J_e\). Lemma~\ref{sci:lem:annihilator},
applied directly to \(I_X\), identifies its nonzero class with the
annihilator line. Compute \(J_e\), \(r\), and the scalar ratio exactly
as above, omitting \((\dagger)\). Fixed \(e\) bounds the quotient
space, while fixed \(p,c\) bounds the number of sparse product choices.
\end{proof}

\begin{example}[A nonreduced regular sublist]
\label{sci:ex:double}
Let \(E(x,y,z)\) be a smooth geometrically integral plane cubic over
\(\Fp\), and work in \(S=\Fp[x,y,z,w]\). The list
\[
 wx,\quad wy,\quad wz,\quad w^2,\quad E
 \]
generates \(K=(\mathfrak m w,E)\), whose saturation is the
anticanonical codimension-two ideal \(I_X=(w,E)\).
It does not contain \(w\). Choose \(f_1=w^2,f_2=E\).
This is a regular sequence with \(e=1\), and every input degree is at
most \(\min(2,3)+1\), so \(B=1\).
The intermediate scheme \(Y\) is nonreduced, and the recovered
annihilator line is generated by \(q=w\), which vanishes on \(X\).
The formula gives
\[
 \mathcal T\!\left(w(w^2E)^{p-1}\right)
 =w\,[(xyz)^{p-1}]E^{p-1}.
 \]
This exhibits a saturation presentation covered by
Theorem~\ref{sci:thm:gap} which does not generate \(I_X\) before
saturation, including the case where a nonvanishing-point
interpretation of \(q\) would be incorrect.
\end{example}

\begin{theorem}[A balanced maximum-degree bound]
\label{sci:thm:balanced}
Fix \(p,c\) and \(E\geq0\). Let a nonzero homogeneous sparse list
\(u_1,\ldots,u_m\) have saturation \(I_X\) with the geometric promises
of Theorem~\ref{sci:thm:gap}. Put \(\delta=\max_i\deg u_i\), and suppose
\[
 0\leq c\delta-(N+1)\leq E.
\]
Then the coherent Frobenius scalar of \(X\) is computable in deterministic
polynomial sparse-bit time, without a designated regular sublist or
an input colon certificate. The algorithm constructs a regular
\(c\)-tuple of linear combinations over an explicitly constructed
finite extension, and descends the scalar to \(\Fp\).
\end{theorem}
\begin{proof}
Replace each \(u_i\) by the \(N+1\) forms
\[
 x_j^{\delta-\deg u_i}u_i\qquad(0\leq j\leq N).
 \]
Write the resulting list as \(U_1,\ldots,U_s\), omitting duplicates
if desired. It has \(s\leq(N+1)m\) and at most \((N+1)T\) terms.
Its saturated ideal is unchanged: on the chart \(x_j\ne0\) the
corresponding multiplier is a unit. All these forms have degree
\(\delta\), whose numerical value is \(O(N+1)\) under the displayed
bound.

Choose a prime
\[
 k>\max(s,\delta^{c-1},p),\qquad D=k^c.
 \]
One may find such a prime of at most twice this threshold by trial
division and Bertrand's bound, in polynomial time in the numerical
threshold. Since \(c,p\) are fixed, \(D\) is polynomial in the input
size. Construct an absolute field \(L=\mathbb F_{p^D}\).
The deterministic fixed-characteristic construction of an irreducible
polynomial of degree \(D\) has polynomial cost in \(D\);
see \citet[Theorem~3.2, author manuscript p.~3]{sci-shoup}.

All nested fields can be found inside this one representation:
\[
 K_r=\{a\in L:a^{p^{k^r}}=a\}
       =\mathbb F_{p^{k^r}},\qquad 0\leq r\leq c.
 \]
Compute their \(\Fp\)-bases as kernels of the corresponding powers
of the Frobenius matrix on \(L\). For \(r<c\), choose any basis
element \(\gamma_r\in K_{r+1}\setminus K_r\). Since
\([K_{r+1}:K_r]=k\) is prime, \(\gamma_r\) has degree exactly \(k\)
over \(K_r\). This avoids any finite-field isomorphism or compatible
tower construction.

Inductively suppose \(f_1,\ldots,f_r\), defined over \(K_r\), are
regular homogeneous forms of degree \(\delta\), each a linear
combination of the \(U_l\). Set
\[
 f_{r+1}=\sum_{l=1}^s\gamma_r^{\,l-1}U_l.
 \]
We check regularity rather than assume a generic choice.
The projective complete intersection cut out by the first \(r\)
forms has degree \(\delta^r\), hence at most \(\delta^r\)
geometric irreducible components. A geometric component has a field
of definition of degree \(d\leq\delta^r<k\) over \(K_r\), because
its Galois orbit has at most that many components. It is geometrically
integral over this field of definition, and its function field is a
regular extension of that finite field. Since \(\gcd(k,d)=1\),
\(\gamma_r\) still has degree \(k\) over the component's function
field. Choose a coordinate \(x_j\) not identically zero on the component
and trivialize \(\OO(\delta)\) on this chart, replacing every section
\(U_l\) by the rational function \(U_l/x_j^\delta\).
The powers \(1,\gamma_r,\ldots,\gamma_r^{s-1}\) are linearly independent
over this function field.

Each such component has dimension \(N-r>N-c=\dim X\), and cannot
lie in the common zero locus \(X\) of all the \(U_l\). At least one
of their restrictions is consequently nonzero. Independence shows
that \(f_{r+1}\) is nonzero on every geometric component. A quotient
by a regular sequence is Cohen--Macaulay and has no embedded
associated primes, so \(f_{r+1}\) is a nonzerodivisor after the field
extension. This proves the induction, including \(r=0\).
Regularity is preserved on extending further to \(L\).

Adjoin the resulting \(f_1,\ldots,f_c\) to the list \(U_l\).
They lie in its ideal, so its saturation remains \(I_{X_L}\).
All forms have degree \(\delta\), and the indexed tuple has excess
\(e=c\delta-(N+1)\leq E\). The linear annihilator construction of
Theorem~\ref{sci:thm:gap} now applies over \(L\) with \(B=0\).
Each \(f_j\) has at most the total number of terms of the \(U_l\);
taking linear combinations has caused no multiplicative increase
in sparse support.

Over \(L\), replace \(\mathcal T\) by the inverse-Frobenius-semilinear
operator which also sends each coefficient \(a\) to \(a^{1/p}\).
The same Laurent pairing proof gives the adjoint Cartier operator.
Find \(q\ne0\) modulo \(J_e\) and compute \(\beta\in L\) with
\(\mathcal C_Y(q)=\beta q\) modulo \(J_e\). The canonical line of
\(X\) is defined over \(\Fp\). If \(\eta\) is its ground-field
generator and the corresponding chosen generator over \(L\) is
\(\lambda\eta\), then
\[
 \beta=\alpha\,\lambda^{1/p}/\lambda,\qquad
 \operatorname{N}_{L/\Fp}(\beta)=\alpha^D.
 \]
This is a norm of a semilinear eigenvalue, not the assertion
\(\beta=\alpha\).
If \(\beta=0\), return zero. Otherwise \(k>p\) implies
\(\gcd(D,p-1)=1\), so exponentiation by the inverse of \(D\)
modulo \(p-1\) recovers \(\alpha\) from the norm. For \(p=2\)
the only nonzero answer is one.

The field dimension \(D\), all subfield linear algebra, coefficient
systems, and term enumerations are polynomial for fixed \(p,c,E\).
Coefficient \(p\)th roots, field norms, and the final exponentiation
also have polynomial bit cost in this explicit field representation.
\end{proof}

\subsection{Finite fibres without a degree gap}

The degree gap bounds the coefficient spaces in \((\dagger)\). The next
algorithm instead evaluates membership in finite algebras. Once a good
projection centre is known, a fixed interpolation grid recovers the
annihilator exactly. Saturation makes this grid independent of the
degrees of the other input equations.

\paragraph{Relation to finite-fibre elimination.}
\citet[Theorem~1, author version p.~6; Definitions~3--4,
pp.~16--17]{sci-gls} use projective Noether position and lifting fibres
to compute geometric resolutions in characteristic zero, assuming reduced
regular prefixes and an invertible fibre Jacobian.
\citet[\S4, author pp.~34--39, especially Propositions~44--46]{sci-hmps}
use Gorenstein duality for interpolation and division in reduced
complete-intersection algebras over number fields. Here the output is
a bounded-degree annihilator line at fixed positive characteristic.
The finite algebras retain their nilpotents, and residual freeness makes
every affine fibre valid.

\begin{lemma}[A common basis for finite fibres]
\label{sci:lem:fibres}
Let \(k\) be a finite field and let \(f_1,\ldots,f_c\) be a homogeneous
regular sequence in \(k[x_0,\ldots,x_N]\), of positive degrees
\(d_1,\ldots,d_c\). Put
\[
 n=N-c\geq1,\qquad \Delta=\prod_jd_j,\qquad
 \rho=\sum_j(d_j-1).
\]
Suppose an \((N+1)\times c\) matrix \(U\) has the property that the
forms \(H_j(z)=f_j(Uz)\) have no common nonzero geometric zero.
Complete its columns to an invertible matrix \([U\ V]\), and write
\(x=Uz+Vt\), where \(t=(t_0,\ldots,t_n)\).
There is a homogeneous monomial list \(\mathcal B\), of length
\(\Delta\) and degrees at most \(\rho\), which is a basis of
\[
 k[x]/(f_1,\ldots,f_c)
 \quad\text{over }k[t_0,\ldots,t_n].
\]
For every \(a\in k^n\), the same list is a \(k\)-basis of
\[
 A_a=k[z_1,\ldots,z_c]/
 \bigl(f_j(Uz+V(1,a)):1\leq j\leq c\bigr).
\]
The basis and multiplication matrices of \(A_a\) can be constructed
by linear algebra in polynomials in \(c\) variables of degrees at most
\(\rho+1\). For fixed \(c\), the cost is polynomial in the numerical
degrees, the displayed sparse size, and \(\log|k|\). This includes
nonreduced fibres.
\end{lemma}
\begin{proof}
The hypothesis implies that \(U\) has rank \(c\): a nonzero vector
in its kernel would be a common zero of all the positive-degree
\(H_j\). The quotient \(k[z]/(H_1,\ldots,H_c)\) is an Artinian
complete intersection. Its Hilbert series is
\[
 \prod_{j=1}^c(1+T+\cdots+T^{d_j-1}),
\]
so it has dimension \(\Delta\) and vanishes in degrees above
\(\rho\). Choose a monomial basis in each degree by coefficient
elimination; their union is \(\mathcal B\).

In the graded ring \(A=k[x]/(f)\), the forms \(t_0,\ldots,t_n\)
are a homogeneous system of parameters. They form a regular sequence
because \(A\) is Cohen--Macaulay. Graded Nakayama gives a surjection
from the graded free \(k[t]\)-module with basis \(\mathcal B\) to
\(A\). Regularity of the parameters gives the same Hilbert series
for both modules, so the kernel is zero in every degree. In particular
specialization at \(t=(1,a)\) gives the asserted basis for every fibre.

To compute its multiplication matrices, put
\(P_j(z)=f_j(Uz+V(1,a))\). Its top \(z\)-homogeneous part is
\(H_j\). For a polynomial of degree \(s\), reduce its degree-\(s\)
part modulo the \(H_j\) to the basis elements of degree \(s\).
Replacing the resulting multiples of \(H_j\) by the same multiples
of \(P_j\) leaves a remainder of smaller degree. Descending in degree
therefore expresses every polynomial of degree at most \(\rho+1\)
as a combination of \(\mathcal B\) and multiples \(a_jP_j\) with
\(\deg a_j\leq\rho+1-d_j\). Apply this to each \(z_i b\),
\(b\in\mathcal B\). The coefficients on \(\mathcal B\) are unique
because it is a fibre basis. Thus one may obtain them by solving the
displayed bounded-degree coefficient systems, without constructing
polynomials in all the parameter variables. The relevant monomial
space has dimension \(\binom{c+\rho+1}{c}\), polynomial for fixed
\(c\). Expansion of the selected forms after substitution takes place
only in these \(c\) variables and has the same polynomial bound.
\end{proof}

\begin{lemma}[A free residual quotient and its fibres]
\label{sci:lem:residual-free}
Work over a finite field \(k\). Let \(J=(f_1,\ldots,f_c)\subseteq I\)
be homogeneous complete-intersection ideals of codimension \(c\) in
\(S=k[x_0,\ldots,x_N]\). Let the degree sums of their regular
generators differ by \(e\geq0\), and let \(U\) be a good centre
for \(f\), as in Lemma~\ref{sci:lem:fibres}. With the same coordinates
\(x=Uz+Vt\), put
\[
 A=k[t_0,\ldots,t_n],\qquad R=S/J,\qquad Q=J:I.
\]
For \(a\in k^n\), write
\(k_a=A/(t_0-1,t_1-a_1,\ldots,t_n-a_n)\) and
\(R_a=R\otimes_A k_a\), and let \(I_a,J_a\)
denote the images of \(I,J\) in \(k[z]\) under this substitution.
Then \(Q/J\) is a graded free direct summand of \(R\) over \(A\),
and \(S/Q\) has a homogeneous free \(A\)-basis in nonnegative
degrees. For every fibre \(t=(1,a)\),
\[
 \operatorname{Ann}_{R_a}(I_a/J_a)
       =(Q/J)\otimes_A k_a.
\]
If \(K^{\mathrm{sat}}=I\), then \(K_a=I_a\) in every such fibre.
These assertions include nonreduced fibres.
\end{lemma}
\begin{proof}
Choose homogeneous regular generators \(g\) of \(I\), write
\(f=Hg\), and put \(q_0=\det H\), of degree \(e\).
Wiebe's identities give
\[
 Q=J+(q_0),\qquad J:q_0=I.
\]
Consequently multiplication by \(q_0\) identifies \(Q/J\) with
\((S/I)(-e)\). The centre \(U\) is good for \(g\) as well:
a common nonzero zero of \(g(Uz)\) would also annihilate \(f(Uz)\).
Lemma~\ref{sci:lem:fibres} therefore makes both \(R\) and
\(M=S/I\) graded free over \(A\).

At \(t=0\), the tuples \(f(Uz)\) and \(g(Uz)\) are homogeneous
regular sequences in \(k[z_1,\ldots,z_c]\). Wiebe's second identity
in this Artinian complete intersection makes the induced map
\[
 M/(t)M\,(-e)\xrightarrow{\ q_0(Uz)\ } R/(t)R
\]
injective. Extend the image basis to a homogeneous basis of \(R/(t)R\)
using monomial classes, and let \(\mathcal D\) be representatives
of the added monomials. Lift a homogeneous monomial basis \(\mathcal B_M\)
of \(M/(t)M\) to \(M\). The elements
\(q_0\mathcal B_M\cup\mathcal D\) generate \(R\) over \(A\)
by graded Nakayama. The free module on these homogeneous generators
has the same Hilbert series as \(R\): modulo \((t)\), their degrees
give exactly the Hilbert series of \(R/(t)R\), and \(R\) is free
over \(A\). The resulting graded surjection is thus an isomorphism.
The first group of basis elements spans \(Q/J\), and the second
gives a free basis of \(S/Q\). Their degrees are nonnegative.

On every affine fibre the two restricted tuples remain regular,
by Lemma~\ref{sci:lem:fibres} applied to both \(f\) and \(g\).
Applying Wiebe's identities to \(f_a=H_a g_a\) gives
\[
 \operatorname{Ann}_{R_a}(I_a/J_a)=q_{0,a}(S/I)_a.
\]
The direct-sum decomposition just proved identifies this with the
base change of \(Q/J\). Finally, saturation gives an \(r\) with
\(\mathfrak m^r I\subseteq K\). Since \(t_0\) is a linear form
in the \(x_i\), one has \(t_0^r I\subseteq K\). Setting \(t_0=1\)
proves \(I_a=K_a\) for every \(a\), as equality of ideals.
\end{proof}

\begin{proposition}[A deterministic grid from a good centre]
\label{sci:prop:fibre-grid}
Fix \(p,c,E\). Under the geometric saturation hypotheses of
Theorem~\ref{sci:thm:gap}, give a
regular sublist \(f_1,\ldots,f_c\) of total degree \(N+1+e\),
where \(0\leq e\leq E\), and a good centre \(U\) over an explicitly
represented finite extension \(k/\Fp\). Then the colon line and
the original coherent Frobenius scalar are computable deterministically
in polynomial sparse-bit time. All other input degrees may be arbitrary
binary integers.

More precisely, put \(n=N-c\), write \(u_1,\ldots,u_s\) for the
original displayed rows of \(K\), and let \(R_a\) be the fibre of
\(S_k/J_k\) at \(t=(1,a)\). The fibre kernel is that of the map
\[
 (S_e)_k\longrightarrow R_a^s,\qquad
 v\longmapsto(v_a u_{1,a},\ldots,v_a u_{s,a}).
\]
After enlarging \(k\) if needed to contain \(e+1\)
distinct elements \(\theta_0,\ldots,\theta_e\), these kernels from the
\(\binom{n+e}{e}\) fibres indexed by
\[
 \Gamma_e=\{(\theta_{\nu_1},\ldots,\theta_{\nu_n}):
          \nu\in\mathbb N^n,\ |\nu|\leq e\}
\]
intersect in exactly \(((J:K)_e)_k\), where \(J=(f_1,\ldots,f_c)\)
and the subscript denotes scalar extension from \(\Fp\) to \(k\).
The kernel-recovery assertion also holds when the tuple is defined
over \(k\). Scalar recovery for such a tuple additionally requires
\(\gcd([k:\Fp],p-1)=1\), and uses the semilinear descent below.
\end{proposition}
\begin{proof}
Put \(Q=J:I_X\). Lemma~\ref{sci:lem:residual-free} gives a
homogeneous free basis \(b_1,\ldots,b_\rho\) of \(S_k/Q_k\) over
\(k[t_0,\ldots,t_n]\), with \(\deg b_j\geq0\), used only
for this proof. A homogeneous \(v\in(S_e)_k\) has a unique image
\[
 \bar v=\sum_j P_j(t)b_j,
 \qquad \deg P_j=e-\deg b_j,
\]
where a negative degree means \(P_j=0\). Thus
\(P_j(1,t_1,\ldots,t_n)\) has total degree at most \(e\).
Dehomogenization at \(t_0=1\) is injective on a fixed homogeneous
degree. On every fibre, the same lemma and \(K_a=I_{X,a}\) show that
\[
 v u_i=0\text{ in }R_a\text{ for every }i
 \quad\Longleftrightarrow\quad
 \bar v=0\text{ in }(S_k/Q_k)\otimes_A k_a.
\]

Evaluation on \(\Gamma_e\) is injective for polynomials of total
degree at most \(e\). Indeed the Newton polynomials
\[
 P_\mu(T)=\prod_{j=1}^n\prod_{r=0}^{\mu_j-1}(T_j-\theta_r),
 \qquad |\mu|\leq e,
\]
form a basis of that polynomial space. Ordered by total degree, their
evaluation matrix on \(\Gamma_e\) is triangular with nonzero diagonal:
an entry vanishes whenever some \(\mu_j>\nu_j\), and the diagonal
is a product of differences of distinct nodes. Consequently the fibre
conditions for all \(a\in\Gamma_e\) hold exactly when every
\(P_j\) above vanishes. Their intersection is therefore
\((Q_e)_k=((J:K)_e)_k\), by Lemma~\ref{sci:lem:annihilator}.

For the algorithm, construct each finite algebra \(R_a\) by
Lemma~\ref{sci:lem:fibres}. Each original coordinate acts by a linear
combination of its \(c\) multiplication matrices and the identity.
Evaluate every sparse \(u_i\) by binary powering in this algebra.
Multiplication by each of the \(M=\binom{N+e}{e}\) degree-\(e\)
monomials gives a matrix with \(M\) columns and at most
\(s\prod_j\deg f_j\) rows. Intersect these kernels over
\(\Gamma_e\). No basis of the unknown residual quotient is computed.

Compute \(J_e\) by bounded-degree multiplication and take the
one-dimensional quotient of the recovered space by \(J_e\).
When the tuple is defined over \(\Fp\), fix a prime-field linear
complement of \(J_e\) in \(S_e\), project to it, and normalize
one nonzero coordinate of a generator to one. The line is defined over
\(\Fp\), so all normalized coordinates lie in \(\Fp\).
This gives \(q\in S_e\setminus J_e\) with \(J:q=I_X\).
Lemma~\ref{sci:lem:trace} computes its scalar.

There are \(\binom{n+e}{e}\) fibres and \(M\) unknown coefficients,
both polynomial for fixed \(E\). Each algebra has dimension
\(\prod_j\deg f_j\leq(N+1+E)^c\). The fixed-\(c\) Macaulay
systems, sparse evaluations by binary powering, kernel intersections,
and final trace product have polynomial bit cost. If \(|k|\leq e\),
then \(k\) has bounded size, and a bounded-degree field extension
providing enough nodes can be constructed by exhaustive irreducibility
testing. Its degree may be chosen congruent to one modulo \(p-1\),
preserving coprimality with \(p-1\) whenever required. Otherwise no
enlargement is required. High input degrees enter
only through their written bit lengths, not through the grid size.
For a tuple defined over an extension, the same kernel proof applies
there; the semilinear norm step is given in
Theorem~\ref{sci:thm:random-excess}.
\end{proof}

\begin{proposition}[A designated regular tuple with unrestricted other degrees]
\label{sci:prop:random-designated}
Fix \(p,c,E\), with \(c\geq2\) and \(E\geq0\). Let
\(K=(u_1,\ldots,u_m)\) have the geometric saturation promise of
Theorem~\ref{sci:thm:gap}. Give an indexed regular sublist
\(f_1,\ldots,f_c\) with
\[
 \sum_j\deg f_j=N+1+e,\qquad 0\leq e\leq E.
\]
For every integer \(a\geq1\), a randomized algorithm of polynomial
bit cost in the input length and \(a\) returns either a correct
\(q\in S_e\setminus J_e\) with \(J:q=I_X\), together with the
original prime-field Frobenius scalar, or failure. Its failure
probability is at most \(2^{-a}\). It never returns an incorrect
answer on the promise. No numerical-degree bound is imposed on any
other input. Repetition gives a Las Vegas algorithm of expected
polynomial bit cost.
\end{proposition}
\begin{proof}
Put \(J=(f_1,\ldots,f_c)\), \(d_j=\deg f_j\),
\(\Delta=\prod_jd_j\), and \(\varepsilon=2^{-a}\).
Construct an extension \(k/\Fp\) with
\[
 |k|>\max\{E+1,4c\Delta/\varepsilon\}.
\]
Its degree is polynomial in the input length and \(a\), independently
of the other input degrees. Use the deterministic fixed-characteristic
irreducible construction of \citet[Theorem~3.2, author manuscript
p.~3]{sci-shoup}. For bounded-bit sampling, index the elements of \(k\)
by their coefficient vectors, sample \(\lceil\log_2|k|\rceil\) bits,
and reduce the resulting integer index modulo \(|k|\). Every atom has
probability at most \(2/|k|\). Induction on the number of independent
coordinates gives a zero-probability bound \(2d/|k|\) for every
nonzero polynomial of total degree at most \(d\).

Sample an \((N+1)\times c\) matrix \(U\). The centre resultant
\[
 \Psi(U)=\operatorname{Res}_z(f_1(Uz),\ldots,f_c(Uz))
\]
is nonzero: a general \(\PP^{c-1}\) misses the projective complete
intersection \(Y\) of dimension \(N-c\). Its degree in the
coefficients of the \(j\)-th form is \(\Delta/d_j\), and those
coefficients have degree \(d_j\) in \(U\). Thus
\(\deg\Psi\leq c\Delta\), and the chance of a bad centre is at
most \(2c\Delta/|k|<\varepsilon\). The resultant degree can also be
read by fixing the other forms: their generic intersection has
\(\Delta/d_j\) points, and vanishing of the remaining form contributes
one linear factor at each point.

The algorithm does not compute \(\Psi\) symbolically. Form the
\(c\)-variable restrictions \(H_j=f_j(Uz)\), put
\(\rho=\sum_j(d_j-1)\), and test whether
\[
 \sum_j H_j k[z]_{\rho+1-d_j}=k[z]_{\rho+1}.
\]
For a good centre, the Artinian Hilbert series of
Lemma~\ref{sci:lem:fibres} proves this equality. Conversely, a common
nonzero geometric zero would annihilate the left side but not every
monomial on the right. The test is exact and has polynomial size for
fixed \(c\). Report failure if it fails; otherwise apply
Proposition~\ref{sci:prop:fibre-grid}. Its deterministic kernels give
the exact colon line and the original scalar, with no further random
choices or exceptional fibres. Every returned answer is correct.
Repeating with \(a=1\) gives expected at most two runs.
\end{proof}

\begin{example}[Why a finite fibre must retain nilpotents]
\label{sci:ex:slices}
Let \(E(x,y,z)\) over \(\Fp\) be a smooth geometrically integral plane cubic with
nonzero \(x^3\)-coefficient. For any binary integer \(M\geq1\), put
\[
 K_M=(w^2,E,wx^M,wy^M,wz^M),\qquad J=(w^2,E).
\]
Its saturation is \((w,E)\): on each of the charts \(x\ne0\),
\(y\ne0\), and \(z\ne0\), an indicated power is a unit, while
the chart \(w\ne0\) contains no zero of \(w^2\).
The regular tuple has excess one, independently of \(M\).
Choose parameters \(z=1,y=b\). Every fibre algebra has basis
\[
 1,x,x^2,w,wx,wx^2,
\]
because it is \(\Fp[x,w]/(E(x,b,1),w^2)\) and the first equation is
a cubic in \(x\). For
\(q=A x+B y+C z+D w\), the equation \(q w=0\) in this algebra is
equivalent to \(A=0\) and \(Bb+C=0\). It is already imposed by
the row \(wz^M\); it then annihilates the other two power rows as
well. One fibre therefore leaves the two-dimensional space spanned by
\(w\) and \(y-bz\). The two fibres \(b=0,1\) leave exactly
\(\Fp w\), the true annihilator line.
Passing to the radical would instead set \(w=0\), making all the
power rows vanish and leaving every degree-one \(q\) admissible.
For instance, over \(\mathbb F_5\),
\(E=y^2z-x^3-xz^2-z^3\) has nine projective points and scalar two;
the recovered class \(w\) gives that scalar for every \(M\).
\end{example}

\subsection{A deterministic algorithm with a linear prefix}

\begin{theorem}[A linear prefix in the actual input ideal]
\label{sci:thm:linear-prefix}
Fix \(p,c,E\), with \(c\geq2\) and \(E\geq0\).
Let \(K=(u_1,\ldots,u_m)\)
satisfy the geometric saturation hypotheses of
Theorem~\ref{sci:thm:gap}. Suppose
\[
 \dim_{\Fp}K_1=c-1,
\]
and, after some finite scalar extension, the actual ideal \(K\)
contains a homogeneous regular \(c\)-sequence of degree sum at most
\(N+1+E\). Then the coherent Frobenius scalar is computable in
deterministic polynomial sparse-bit time. No regular tuple or projection
centre is supplied, the linear forms may involve every coordinate,
and the other input degrees may be arbitrary binary integers.
\end{theorem}
\begin{proof}
The displayed degree-one inputs span \(K_1\). Gaussian elimination
gives independent forms \(L_1,\ldots,L_{c-1}\) over \(\Fp\).
Put
\[
 P=(L_1,\ldots,L_{c-1}),\qquad W=V(P)\subset\mathbb A^{N+1},
 \qquad \delta=N-c+2,\qquad D=\delta+E.
\]
Some displayed input of degree at most \(D\) lies outside \(P\).
Indeed a regular \(c\)-sequence cannot be contained in the
height-\((c-1)\) ideal \(P\). Choose one of its members outside
\(P\). Its degree is at most \(N+1+E-(c-1)=D\), because the
other \(c-1\) degrees are positive. Its homogeneous expression in
the inputs uses only \(u_i\) of degree at most \(D\), and at least
one of these must lie outside \(P\). The argument remains valid after
scalar extension, since membership in \(P\) commutes with it.

We find such an input without expanding its restriction to \(W\).
Row reduction selects \(c-1\) pivot coordinates and expresses them
as linear forms \(M_1(y),\ldots,M_{c-1}(y)\) in the remaining
\(\delta\) coordinates \(y\). Each restricted input monomial is
of the form
\[
 a y^\alpha\prod_{j=1}^{c-1}M_j(y)^{b_j}.
\]
For an input of degree at most \(D\), this is a semidiagonal circuit
with at most \(c-1\) noncoordinate linear powers per product gate
and formal degree at most \(D\). Its top fan-in is at most the
original term count. The deterministic identity test of
\citet[Definition~1.1, Corollary~3.3, and \S6, especially
Theorems~6.3 and~6.5]{sci-sss} applies in every characteristic and
also permits affine linear factors. Scan these low-degree inputs and
select one with nonzero restriction, denoted \(f\).

The source's small-characteristic algorithm has polynomial bit cost
in the present setting. For \(k=\Fp[T]/(\mu)\) of degree \(h\),
its lift uses
\((\mathbb Z/p^{B+1}\mathbb Z)[T]/(\widetilde\mu)\), with
\(\widetilde\mu\) the monic coefficientwise integer lift and
\(B\) polynomial in the numerical degrees and number of variables.
Its elements have \(h\) coefficients of \(O(B\log p)\) bits;
arithmetic and inversion of elements nonzero modulo \(p\) are
polynomial-time operations. The module-pruning algorithm keeps
polynomially many coefficient rows. Since \(c-1\) is fixed, the
resulting semidiagonal test is polynomial in the sparse size, \(D\),
and \(h\). No unrestricted-degree input is submitted to this test.

The tuple \(L_1,\ldots,L_{c-1},f\) is regular, since \(S/P\)
is a polynomial domain and the image of \(f\) is nonzero. Its degree
\(d=\deg f\) satisfies \(\delta\leq d\leq\delta+E\).
For the lower bound, write the tuple in terms of anticanonical regular
generators of \(I_X\). Wiebe's identity makes the coefficient
determinant nonzero; its degree is \((c-1)+d-(N+1)=d-\delta\),
which cannot be negative. Thus this tuple has excess
\(e=d-\delta\leq E\).

Construct \(k/\Fp\) with \(|k|>D\), taking the least suitable
power of \(p\). Its cardinality is at most \(p(D+1)\), so exhaustive
irreducible construction has polynomial cost at fixed \(p\).
Find \(a\in W(k)\) with \(f(a)\ne0\) by specializing the free
coordinates one at a time. At each stage try \(D+1\) distinct field
values and retain one leaving a nonzero semidiagonal polynomial.
At most \(D\) choices can annihilate the current nonzero polynomial:
any one nonzero coefficient, viewed as a polynomial in the coordinate
being specialized, has degree at most \(D\). Partial specialization
keeps the monomial factor and at most \(c-1\) affine linear powers,
so the same identity test applies. This needs at most
\(\delta(D+1)\) tests. Homogeneity and \(f(a)\ne0\) imply \(a\ne0\).

Choose vectors \(v_1,\ldots,v_{c-1}\) over \(\Fp\) with
\(L_i(v_j)=\delta_{ij}\), using linear algebra, and let
\[
 U=[v_1\ \cdots\ v_{c-1}\ a].
\]
Its columns are independent. The restrictions of the first \(c-1\)
equations are \(L_i(Uz)=z_i\). At their common zero the last is
\(f(a)z_c^d\), which has no nonzero zero. Hence \(U\) is a good
centre, constructed deterministically.

Adjoin the \(L_i\) to the input list, without changing its ideal,
and apply Proposition~\ref{sci:prop:fibre-grid} to this tuple and
centre. The tuple is defined over \(\Fp\), so normalization descends
the recovered colon line to \(\Fp\), even though the centre uses
\(k\). The original sparse forms suffice for the Cartier calculation.
All candidate degrees are numerically \(O(N+E)\); all other inputs
are used only by binary powering in polynomial-dimensional fibre
algebras. The preceding bounds prove deterministic polynomial bit cost.
\end{proof}

\subsection{A deterministic algorithm with a factored linear prefix}

\begin{theorem}[A supplied factored prefix]
\label{sci:thm:factored-prefix}
Fix \(p,c,E\), with \(c\geq2\) and \(E\geq0\), and let \(K\)
satisfy the geometric saturation hypotheses of
Theorem~\ref{sci:thm:gap}. Suppose distinct supplied indices designate
original input rows \(f_1,\ldots,f_c\) forming a homogeneous regular
sequence with
\[
 \sum_{i=1}^c\deg f_i=N+1+e,\qquad 0\leq e\leq E.
\]
For each \(i<c\), also supply a promised exact factorization
\[
 f_i=a_i\prod_{t=1}^{s_i}L_{it}^{m_{it}},
 \qquad a_i\in\Fp^*,\quad m_{it}\geq1,
\]
where the nonzero homogeneous linear forms \(L_{it}\) are given by
their coefficient vectors over \(\Fp\), and the positive multiplicities
are given in binary. Repeated factors are permitted. Then the coherent
Frobenius scalar is computable in deterministic polynomial sparse-bit
time. No centre, colon generator, or linear form in \(K\) is supplied,
and all other input degrees may be arbitrary binary integers.
No factorization algorithm or verification of the promised factorization
equalities is assumed.
\end{theorem}
\begin{proof}
Put \(r=c-1\), \(R=N+1+E\), and \(d=\deg f_c\). Choose one
listed factor from each of \(f_1,\ldots,f_r\), and put their coefficient
vectors into the rows of a matrix \(A\) of size \(r\times(N+1)\).
There are
\[
 T=\prod_{i=1}^r s_i\leq\prod_{i=1}^r\deg f_i\leq R^r
\]
such choices; duplicate matrices or subspaces may be retained.
Every \(A\) has rank \(r\). Indeed its linear ideal contains
\((f_1,\ldots,f_r)\), which has height \(r\), so
\(r\leq\operatorname{rank}A\leq r\).
The spaces \(W_A=\ker A\) cover the geometric zero set of this prefix.
Moreover
\[
 f_c|_{W_A}\ne0\qquad\hbox{for every }A.
 \tag{\(\ddagger\)}
\]
Otherwise the full tuple would lie in the height-\(r\) prime ideal
generated by the selected linear forms, contradicting its regularity.
These statements concern the components of the supplied prefix; none
of its individual factors is asserted to belong to \(K\).

Choose an explicit extension \(k/\Fp\) whose cardinality is the least
power of \(p\) exceeding
\[
 T(N+d)+d+E+2.
\]
Its cardinality is polynomial in the written input size at fixed
\(p,c,E\), and its extension degree is logarithmic in that bound.
Enumerating monic polynomials of this degree and applying deterministic
irreducibility testing constructs \(k\) in polynomial time. This one
field supplies every search below and the final degree-\(e\) grid.
No coprimality restriction on \([k:\Fp]\) is necessary, because the
original tuple remains defined over \(\Fp\).

First construct an \((N+1)\times r\) matrix \(B\) for which every
\(AB\) is invertible. Suppose \(j<r\) columns have been chosen and
their image under each \(A\) has rank \(j\). For every \(A\),
compute a nonzero row \(\lambda_A\) annihilating those image columns.
Since \(A\) has full row rank, \(\lambda_AA\ne0\). On the vector
\[
 b(t)=(1,t,\ldots,t^N)^{\mathsf t},
\]
the polynomial \(\lambda_A A b(t)\) is nonzero and has degree at most
\(N\). The union of its zeros over all \(T\) choices has at most
\(TN\) elements. Testing \(TN+1\) distinct field values finds a
column outside every forbidden span. Repeating this \(r\) times gives
\[
 AB\in\operatorname{GL}_r(k)\quad\hbox{for every }A.
\]
Only explicit linear algebra and univariate evaluations are used.

Next find, separately for every \(A\), a point
\[
 a_A\in W_A(k),\qquad f_c(a_A)\ne0.
\]
Row reduction expresses \(r\) pivot coordinates as linear forms
\(M_1(y),\ldots,M_r(y)\) in the other \(N-c+2\) coordinates.
Substituting into a sparse monomial of \(f_c\) gives
\[
 \gamma y^\alpha\prod_{j=1}^r M_j(y)^{b_j}.
\]
The restriction is therefore a semidiagonal circuit with at most the
original number of terms, at most \(r\) noncoordinate linear powers
per term, and formal degree at most \(d\leq R\). It is nonzero by
\((\ddagger)\). Apply the characteristic-independent test of
\citet[Definition~1.1, Corollary~3.3, and \S6,
Theorems~6.3 and~6.5]{sci-sss}. Specialize each free coordinate in
turn, trying \(d+1\) distinct values and retaining one for which the
remaining polynomial is nonzero. At most \(d\) values annihilate
any one nonzero coefficient polynomial in that coordinate. Partial
specialization preserves the semidiagonal form, with affine linear
factors allowed. This gives \(a_A\) in at most
\((N-c+2)(d+1)\) tests.

The small-characteristic bit implementation is the one detailed in
Theorem~\ref{sci:thm:linear-prefix}: over an explicit degree-\(h\)
finite extension, the SSS lift uses a Galois ring represented by \(h\)
coefficients modulo \(p^{B_0+1}\), with \(B_0\) polynomial in the
numerical degree and variable bounds. Its duality and module-pruning
steps have polynomial cost for the fixed parameter \(r\).
Here \(h=O(\log(T(N+d)+d+E+2))\) and \(d\leq R\), so this is a
polynomial bit bound at every fixed prime. None of the other,
unrestricted-degree rows is used in these identity tests.

For every \(A\), form the explicit projection
\[
 P_A=I_{N+1}-B(AB)^{-1}A.
\]
It maps onto \(W_A\), fixes \(W_A\), and kills the column span of
\(B\). We combine the separate points into one vector \(v\) such
that
\[
 f_c(P_Av)\ne0\qquad\hbox{for all }A.
 \tag{\(\clubsuit\)}
\]
Order the matrices as \(A_1,\ldots,A_T\) and start with \(v=0\).
At stage \(j\), replace \(v\) by \(v+t a_{A_j}\).
For each already satisfied condition \(i<j\), the polynomial
\[
 f_c(P_{A_i}v+tP_{A_i}a_{A_j})
\]
is nonzero because its value at \(t=0\) is nonzero. For the new
condition, the coefficient of \(t^d\) in
\[
 f_c(P_{A_j}v+t a_{A_j})
\]
is \(f_c(a_{A_j})\ne0\), by homogeneity. Each polynomial has degree
at most \(d\), so the \(j\) conditions together forbid at most
\(jd\leq Td\) values. Test \(Td+1\) values and keep one satisfying
them all. These tests evaluate the original sparse \(f_c\) at explicit
vectors; they do not expand any projected polynomial. After \(T\)
stages, \((\clubsuit)\) holds. All loops have polynomial bounds.

Set \(U=[B\ v]\). Its columns are independent: otherwise
\(P_Av=0\) for every \(A\), contradicting \((\clubsuit)\) and
\(d>0\). At a common geometric zero \(z=(z',z_c)\) of the restricted
tuple, each product \(f_i(Uz)\), \(i<c\), supplies a vanishing
listed linear factor. The resulting matrix \(A\) gives
\[
 ABz'+Avz_c=0,
 \qquad Uz=z_cP_Av.
\]
Consequently
\(f_c(Uz)=z_c^d f_c(P_Av)=0\) forces \(z_c=0\), and then
\(z'=0\). Thus \(U\) is a good centre over the algebraic closure.
Repeated factors have only been ignored for this set-theoretic centre
argument; all original equations and multiplicities are retained below.

Apply Proposition~\ref{sci:prop:fibre-grid} to the original designated
prime-field tuple and the full original input list. Its colon line
descends to \(\Fp\) by the prime-field normalization there. Equivalently,
expand the extension-field kernel equations in a fixed \(\Fp\)-basis
of \(k\), restrict the unknown coefficients to \(\Fp\), and select a
class outside \(J_e\). The ensuing Cartier calculation therefore
returns the original prime-field scalar without norm ambiguity.
The finite algebras have dimension at most \(R^c\), the degree-\(e\)
unknown space and fibre grid are polynomial for fixed \(E\), and the
trace uses the original sparse tuple with fixed \(p,c\). All other
rows enter only through binary powering in these finite algebras.
This proves the full polynomial sparse-bit bound.
\end{proof}

\begin{example}[A split prefix without actual linear generators]
\label{sci:ex:factored-prefix}
Let \(I_X=(w,G(x_1,\ldots,x_N))\), where \(N\geq3\) and \(G\)
is a sparse equation of a smooth geometrically integral degree-\(N\)
hypersurface in \(\PP^{N-1}\). For a binary integer \(M\geq2\), put
\[
 K=(w(w+x_1),\ G,\ wx_1^M,\ldots,wx_N^M).
\]
On every chart \(x_i\ne0\), a cover row gives \(w\), so the localized
ideal equals \(I_X\). On \(w\ne0\), the equations \(w+x_1=0\)
and \(x_1^M=0\) give the unit ideal. Hence \(K^{\mathrm{sat}}=I_X\).
The designated pair \((w(w+x_1),G)\) is regular: \(G\) is divisible
by neither \(w\) nor \(w+x_1\). Its excess is one and its first
factorization is supplied explicitly. Yet \(K_1=0\), so the actual
linear-prefix hypothesis does not hold. Theorem~\ref{sci:thm:factored-prefix}
applies with unrestricted binary \(M\).
\end{example}

\subsection{Discovering a bounded-excess sequence}

\begin{lemma}[Deterministic candidates for unequal degree profiles]
\label{sci:lem:profile-candidates}
Fix \(p,c,E\), and let \(K\) satisfy the geometric saturation
hypotheses above. Put \(R=N+1+E\). For every nondecreasing positive
profile \(d_1\leq\cdots\leq d_c\) with sum in \([N+1,R]\),
one can deterministically construct an explicitly sparse candidate tuple
of the prescribed degrees over a common finite extension of polynomial
degree. If \(K_{\leq d}\) denotes the ideal generated by the displayed
inputs of degree at most \(d\), then every profile satisfying
\[
 \operatorname{ht}K_{\leq d_j}\geq j\qquad(1\leq j\leq c)
 \tag{\(\star\)}
\]
produces a regular tuple. The extension degree may be chosen coprime to
\(p-1\). If \(K\), after scalar extension, contains a regular tuple
of degree sum at most \(R\), at least one candidate is regular.
The lemma does not assert recognition of the successful candidates.
\end{lemma}
\begin{proof}
Every regular \(c\)-tuple in \(K\) has degree sum at least \(N+1\).
Indeed write it in terms of anticanonical regular generators of \(I_X\).
Wiebe's identity makes the determinant of the coefficient matrix nonzero,
and its degree is the difference of the two degree sums. A homogeneous
polynomial of negative degree is zero, so that difference is nonnegative.
One may sort the degrees: regular sequences in the polynomial ring
remain regular after permutation, since each prefix has the required
height and the ring is Cohen--Macaulay. A regular tuple of sorted
profile \(d\) implies \((\star)\), because its first \(j\) forms
use only generators of degrees at most \(d_j\).

For each profile form the sparse lists
\[
 \mathcal W_j=
 \{x_l^{d_j-\deg u_i}u_i:
       \deg u_i\leq d_j,\ 0\leq l\leq N\}.
\]
Their projective base loci equal those of \(K_{\leq d_j}\): at a
projective point at least one coordinate is nonzero. Each list has
length at most \((N+1)m\) and total sparse size at most \((N+1)\)
times that of the original input.

Choose an integer prime
\[
 \ell>\max\{(N+1)m,R^{c-1},p\},\qquad h=\ell^c,
\]
and work in \(L=\mathbb F_{p^h}\). As in
Theorem~\ref{sci:thm:balanced}, construct this absolute field
deterministically, realize its nested subfields
\(k_r=\mathbb F_{p^{\ell^r}}\) by Frobenius kernels, and choose
\(\gamma_r\in k_{r+1}\setminus k_r\). Every \(\gamma_r\) has
degree \(\ell\) over \(k_r\). Enumerating
\(\mathcal W_j=(W_{j1},\ldots,W_{js_j})\), define
\[
 f_j=\sum_{l=1}^{s_j}\gamma_{j-1}^{\,l-1}W_{jl}.
\]
An empty list gives a zero candidate. All candidates are constructed
without testing \((\star)\).

Suppose now that \((\star)\) holds. Inductively, a regular prefix
of length \(j-1\), defined over \(k_{j-1}\), has degree
\(\prod_{r<j}d_r\leq R^{j-1}<\ell\). Each geometric component
has field of definition of degree less than \(\ell\) over
\(k_{j-1}\), since its Galois orbit has no more elements than the
number of components. Its function field is a regular extension of
that finite field. The prime degree \(\ell\) therefore remains
the degree of \(\gamma_{j-1}\) over the component's function field.
The component has codimension \(j-1\), whereas the base locus of
\(\mathcal W_j\) has codimension at least \(j\). Some section
in \(\mathcal W_j\) is nonzero on it. Trivialize these degree-\(d_j\)
sections by dividing by \(x_v^{d_j}\) on a nonempty coordinate chart.
Since \(s_j<\ell\), independence of
\(1,\gamma_{j-1},\ldots,\gamma_{j-1}^{s_j-1}\) over the function
field proves that \(f_j\) is nonzero on this component. The preceding
complete intersection is Cohen--Macaulay with no embedded associated
primes, so \(f_j\) is a nonzerodivisor. This proves regularity.
Here every relevant cone component has nonempty projectivization,
since \(j\leq c\leq N-1\).

There are at most \(R^c\) profiles. Bertrand's bound and trial
division find a suitable \(\ell\) of polynomial size, and any larger
polynomially bounded threshold may be imposed on \(\ell\).
The absolute extension degree \(h=\ell^c\), all subfield linear
algebra, and all sparse combinations have polynomial size for fixed
\(p,c,E\). Since \(\ell>p\), one has \(\gcd(h,p-1)=1\).
The first paragraph shows that a promised regular tuple contributes
at least one profile satisfying \((\star)\), completing the claim.
\end{proof}

\begin{theorem}[Bounded complete-intersection excess in the input ideal]
\label{sci:thm:random-excess}
Fix \(p,c,E\), with \(c\geq2\) and \(E\geq0\). Let a homogeneous
sparse ideal \(K=(u_1,\ldots,u_m)\) have the geometric saturation
promise of Theorem~\ref{sci:thm:gap}. Suppose that, after some finite
scalar extension, the actual ideal \(K\) contains a homogeneous
regular sequence of length \(c\) whose degree sum is at most
\(N+1+E\). Neither the sequence, its degrees, nor the extension is
supplied. Equivalently, such a sequence may be promised over
\(\overline{\Fp}\).

For every integer \(a\geq1\), a randomized algorithm of polynomial bit cost
in the sparse input length and \(a\) returns the original
\(\Fp\)-coherent Frobenius scalar or reports failure, with failure
probability at most \(2^{-a}\). Every returned scalar is correct on
the promise. It follows that scalar recovery has a Las Vegas algorithm
of expected polynomial bit cost. All original input degrees may be
arbitrary binary integers; the bounded-excess condition is imposed on
a sequence in \(K\), not merely in \(K^{\mathrm{sat}}\).
\end{theorem}
\begin{proof}
Put \(R=N+1+E\) and \(\varepsilon=2^{-a}\). Choose an integer prime
\[
 \ell>\max\!\left\{(N+1)m,R^{c-1},p,
 \left\lceil\log_p(4cR^c/\varepsilon+E+2)\right\rceil\right\},
 \qquad h=\ell^c,
\]
and construct \(L=\mathbb F_{p^h}\). This has polynomial degree in
the sparse input length and \(a\), satisfies
\[
 |L|>\max\{E+1,4cR^c/\varepsilon\},
 \qquad \gcd(h,p-1)=1,
\]
and supports all deterministic candidates of
Lemma~\ref{sci:lem:profile-candidates}. Construct those candidates,
one for every sorted positive profile of sum in \([N+1,R]\).
The promise guarantees that at least one is regular. All their displayed
forms belong to the extended actual ideal \(K_L\).

For each candidate, sample only an \((N+1)\times c\) matrix \(U\),
using the bounded-bit sampler of
Proposition~\ref{sci:prop:random-designated}. Apply the exact
Macaulay span test there to \(f_1(Uz),\ldots,f_c(Uz)\), in degree
\(\rho+1=\sum_j(d_j-1)+1\). A passing test certifies global
regularity as well as a good centre. Otherwise the common projective
zero locus of the candidate tuple would have dimension at least
\(N-c+1\) and would meet every \(\PP^{c-1}\), contradicting the
test. Equivalently the global ideal has height \(c\), so its \(c\)
generators form a regular sequence. This certifies every accepted
candidate, even when its profile did not satisfy \((\star)\).

For any one regular candidate, its centre resultant is a nonzero
polynomial of degree at most \(c\prod_jd_j\leq cR^c\). The probability
that this candidate's sampled centre fails is at most
\(2cR^c/|L|<\varepsilon\). If an earlier candidate passes, it is
already usable. Thus no union bound over the profiles is needed.
If every test fails, report failure. The coefficients of the candidate
tuples are deterministic; only their centres are sampled.

On accepting a tuple and centre, put
\(J=(f_1,\ldots,f_c)\) and \(e=\sum_jd_j-(N+1)\).
Adjoin the tuple to the input list, which does not change its ideal.
Apply Proposition~\ref{sci:prop:fibre-grid} over \(L\) to recover
\((J:K_L)_e\) exactly. Its quotient by \(J_e\) is the annihilator
line, and any nonzero class \(q\) on it satisfies \(J:q=I_{X_L}\).
This step is deterministic and has no exceptional fibres, independently
of the other input degrees.

The newly constructed \(J\) and its displayed colon line need not
be defined over \(\Fp\). Use the inverse-Frobenius-semilinear trace
operator over \(L\), taking \(p\)th roots of coefficients as in
Theorem~\ref{sci:thm:balanced}, and read
\(\mathcal C_Y(q)=\beta q\) modulo \(J_e\).
The line is the image of the canonical line of \(X_L\), whose
prime-field generator has scalar \(\alpha\). If \(q\) is
\(\lambda\) times that generator's image, then
\[
 \beta=\alpha\lambda^{1/p}/\lambda,
 \qquad \operatorname{N}_{L/\Fp}(\beta)=\alpha^h.
\]
If \(\beta=0\), return zero. Otherwise, for \(p>2\), choose
\(v\) with \(hv\equiv1\pmod{p-1}\), and return
\(\operatorname{N}_{L/\Fp}(\beta)^v=\alpha\).
For \(p=2\), return one in the nonzero case. This is exact descent
of a semilinear eigenvalue; it does not require the chosen colon line
itself to descend in its displayed coordinates.

The candidate lemma bounds the profiles, sparse lists, and field
representation polynomially. Every fixed-\(c\) Macaulay space is
polynomial in \(R\); every degree-\(e\) coefficient space and
interpolation grid is polynomial for fixed \(E\). High-degree original
inputs are evaluated only by binary powering in the finite fibre
algebras. They are never expanded after the full coordinate change,
and their numerical degrees do not enter the centre probability bound.
The final trace product has at most
\(\binom{N+e}{e}T'^{c(p-1)}\) ordered term choices, where \(T'\)
is the polynomially bounded total term count of the constructed tuple.
Coefficient \(p\)th roots and the norm have polynomial bit cost in
\(h\). Every successful output is correct, and repeating with \(a=1\)
gives expected at most two runs.
\end{proof}

\subsection{Safe degree reduction before saturation recovery}

\begin{lemma}[Digit deflation for a radical saturation]
\label{sci:lem:deflation}
Let a homogeneous sparse ideal \(K\) over \(\Fp\) have radical
saturated ideal \(I=K^{\mathrm{sat}}\), and discard any zero rows.
For \(P=p^r\) with \(r\geq1\), split each exponent vector of one
remaining row as \(a+Pb\), with \(0\leq a_j<P\).
If the sparse coefficient matrix indexed by \(a,b\) has rank one,
write that row as \(UV^P\) and replace it by \(UV\), whenever
\(V\) is nonconstant. Repeated application preserves the saturation,
never increases term counts, and has deterministic polynomial sparse-bit
cost. Whole Frobenius-power extraction is included. No individual
monomial division is asserted.
\end{lemma}
\begin{proof}
The rank-one reconstruction and homogeneity argument are those of
Lemma~\ref{sp:lem:preprocessing}. Its nonzero support is Cartesian,
so the number of term pairs in \(UV\) equals the old term count.
Radicality gives
\[
 (UV)^P=U^{P-1}(UV^P)\in I\quad\Longrightarrow\quad UV\in I.
\]
The old row is a multiple of the new one. Thus every replacement
satisfies \(K\subseteq K'\subseteq I\), and \((K')^{\mathrm{sat}}=I\).
If the row uses \(v\) coordinates, put \(a=\deg U\) and
\(b=\deg V\geq1\). Then \(a\leq v(P-1)\), and its old degree
obeys \(a+Pb\leq(v+2)(P-1)b\). The decrease \((P-1)b\) is
therefore at least the fraction \(1/(v+2)\) of the old degree.
There are \(O(v\log(D+1))\) nontrivial steps from initial degree
\(D\). Each step tests only \(O(\log D)\) digit scales by sparse
rank-one tests, and no exponent bit length increases. This proves
the complexity claim, with the largest scales tested first if a
deterministic ordering is desired.
\end{proof}

For example, in the coordinate-hyperplane presentation of
Example~\ref{sci:ex:double}, dividing the monomial row \(w^2\) by
its entire monomial content would give \(1\notin I_X\).

The hypotheses of Theorems~\ref{sci:thm:random-excess}
and~\ref{sci:thm:linear-prefix} may consequently be imposed on the
actual ideal after this preprocessing. This genuinely permits
unbounded complete-intersection excess in the original ideal.
For example, suppose \(I_X=(g_1,\ldots,g_c)\) has sparse regular
generators of degrees \(d_1\leq\cdots\leq d_c\), and take the
full list
\[
 \{x_j^{P}g_i:1\leq i\leq c,\ 0\leq j\leq N\},
 \qquad P=p^r>\max_i d_i.
\]
Its saturation is \(I_X\), whereas every nonzero element of the
actual ideal has degree at least \(P+\min_i d_i\).
At digit scale \(P\), each row becomes \(x_jg_i\). For the profile
\(d_i+1\), the first \(i\) groups have projective base locus
\(V(g_1,\ldots,g_i)\). The height criterion in
Lemma~\ref{sci:lem:profile-candidates} gives a regular tuple with
degree sum \(N+1+c\). Thus Theorem~\ref{sci:thm:random-excess}
applies after deflation with \(E=c\), regardless of \(r\).

There is also a deterministic example. Suppose
\(I_X=(L_1,\ldots,L_{c-1},F)\), where the \(L_i\) are independent
linear forms and \(\deg F=\delta=N-c+2\), with the required
smooth geometric properties. For \(P=p^r>\delta\), put
\[
 K=(F,\ L_i^P,\ x_j^P L_i:
          1\leq i<c,\ 0\leq j\leq N).
\]
The cover rows on every coordinate chart prove \(K^{\mathrm{sat}}=I_X\).
Whole-power extraction introduces the \(L_i\) into the actual ideal,
so Theorem~\ref{sci:thm:linear-prefix} applies with excess zero.
Originally all elements of degree less than \(P\) are multiples of
\(F\). Any sorted regular \(c\)-profile therefore has its last
\(c-1\) degrees at least \(P\), and its total excess is at least
\((c-1)(P-1)\). The cover rows are essential: the pure-power rows
and \(F\) alone generally define a nonreduced saturation.

\paragraph{Limits of the results.}
Theorem~\ref{sci:thm:random-excess} removes the indexed tuple and the
all-input degree gap under a bounded-excess promise on the actual
input ideal after scalar extension. Smooth anticanonical saturation
alone does not supply that promise. If \(I_X=(g_1,\ldots,g_c)\), the
full list \(x_j^M g_i\), for every \(0\leq j\leq N\) and
\(1\leq i\leq c\), has saturation \(I_X\): on each coordinate
chart the corresponding multiplier is a unit. Its least nonzero degree
is at least \(M+\min_i\deg g_i\), so the least complete-intersection
excess in the actual input ideal can grow without bound.
The deterministic bounded-gap, balanced-maximum-degree, linear-prefix,
and supplied factored-prefix theorems give stronger algorithmic
conclusions on their respective subclasses. Without the gap, the direct coefficient system
\((\dagger)\) may have large coefficient spaces; the deterministic
fibre grid avoids those spaces without expanding them. The unequal-profile
candidate lemma also removes random coefficient selection. In the
bounded-excess range, the remaining derandomization problem is therefore
to find a certified good centre for at least one candidate tuple.
Providing such a centre gives a deterministic algorithm by
Proposition~\ref{sci:prop:fibre-grid}. Without bounded
excess, the degree-\(e\) annihilator calculation itself may be
exponentially large. The explicit-certificate corollary is deterministic
and does not require kernel discovery. None of these results settles
unrestricted fixed-codimension saturation. The following results further
enlarge the tractable class by prime-field relations and certified
coordinate division; the
remaining question concerns unbounded excess after all the stated
safe preprocessing.
